\section{跨模态迁移：先把世界表示成可通信的元素}

模型家族回答“token 之间能怎样连接”，跨模态应用进一步回答“什么算 token”。Karpathy 用一句故意夸张的话概括迁移 recipe：把图像、语音、轨迹或分子结构切成元素，假装它们是文本，再让 Transformer 建模关系。真正的工程难点恰好藏在“切成元素”这一步。

\subsection{Vision Transformer：patch 是视觉 token}

上一章已经说明 encoder 能让一组 token 双向通信；视觉迁移的第一步，就是定义怎样把二维像素变成这组 token。ViT 选择最直接的接口：把图像划分为固定 patch，每个 patch 展平并线性投影为向量，再加入位置表示送入 encoder。若图像大小为 $H\times W$、patch 边长为 $P$，token 数约为 $N=HW/P^2$。更小 patch 提供更细粒度，但使 attention 的 $N^2$ 计算与显存迅速增长，因此 tokenization 同时决定视觉分辨率和系统成本。

\lecturefigure{slide-47-vision-transformer.jpg}{Vision Transformer：把图像切成 patch sequence，再用 encoder 全局通信。}{Stanford Online 官方视频 00:52:41--00:53:42。}

\begin{knowledgebox}{读图：ViT 没有让二维结构消失}
Patch ordering 与 positional embedding 仍编码空间位置，训练增强与数据规模也提供视觉 bias。纯 Transformer 需要重新学习局部性与平移规律，因此常比 ConvNet 更依赖数据或预训练；但全局 attention 让远距离 patch 在浅层即可直接通信，并提供统一接口。
\end{knowledgebox}

\subsection{Speech：spectrogram 序列与局部卷积 bias}

语音可以先变成 time--frequency spectrogram，再沿时间切成帧或块。Conformer 不是“只把语谱图扔给 vanilla Transformer”，而是在 attention 之外加入 convolution module，以同时建模全局依赖与局部声学模式；Whisper 也使用专门前端、encoder--decoder 与大规模弱监督目标。

\lecturefigure{slide-48-conformer.jpg}{Conformer：attention 的全局建模与 convolution 的局部声学 bias 结合。}{Stanford Online 官方视频 00:53:42--00:54:00。}

\begin{knowledgebox}{读图：跨领域复用不是逐层完全相同}
Mel 语谱图是语音在时间与频率上的能量表示；局部卷积擅长捕捉邻近帧模式，自注意力擅长远距离关系。Conformer 的贡献正是把二者组合。课堂所说“把它当成文本”应理解为共享序列接口，而不是抹掉声学结构与专门目标函数。
\end{knowledgebox}

\subsection{Decision Transformer：把控制轨迹写成序列}

Decision Transformer 把 return-to-go、state 与 action 交错排列，让 causal model 在期望回报和历史条件下预测 action。它把 offline RL 重写为 conditional sequence modeling：训练数据来自已有轨迹，不在线与环境交互；规划能力依赖数据覆盖、return conditioning 与模型能否从序列统计中提取行为策略。

\lecturefigure{slide-49-decision-transformer.jpg}{Decision Transformer：return、state、action 组成可自回归建模的轨迹序列。}{Stanford Online 官方视频 00:54:00--00:54:13。}

\begin{knowledgebox}{读图：它不是把奖励当普通单词}
不同元素有不同语义与 embedding，时间步和 modality/type 信息必须被编码。模型学习 $p(a_t\mid R_t,s_{\le t},a_{<t})$ 一类条件分布，而不是直接求解 Bellman equation。序列建模提供统一训练接口，但不消除 offline data bias 与 extrapolation 风险。
\end{knowledgebox}

\subsection{AlphaFold2：attention 进入专门的生物结构系统}

AlphaFold2 的 Evoformer 在 multiple-sequence alignment 与 residue-pair representation 之间执行多种 attention 和 specialized updates，之后结构模块输出三维构象。它说明 attention 可处理“哪些残基与哪些残基相关”，却不能被缩写为 GPT 对氨基酸做普通 next-token prediction。

\lecturefigure{slide-50-alphafold.jpg}{AlphaFold2：attention 位于高度领域化的蛋白质结构推断系统中。}{Stanford Online 官方视频 00:54:13--00:54:23。}

\begin{warningbox}{读图时容易犯的错误：把 backbone 相似写成系统等价}
ViT、Conformer、Decision Transformer 与 AlphaFold2 都使用 attention，但输入单位、position geometry、目标函数、模块组合和输出语义不同。课程想证明的是“可寻址通信原语具有普适性”，不是“所有问题都已经被语言建模统一”。
\end{warningbox}

\subsection{异构条件输入：Transformer 为什么容易“加一种传感器”}

卷积网络的 feature map 与二维几何紧密绑定，添加 radar、map、vehicle state 或 audio 时必须决定在哪一层 concatenate 或 fuse。Transformer 可以把每种输入编码成统一 width 的元素，加入 type embedding 区分来源，再放进允许通信的集合；attention 学习数据依赖的融合权重。

\lecturefigure{slide-51-tesla-input-flexibility.jpg}{Tesla 经验例：异构信息先编码成元素，再交给 self-attention 融合。}{Stanford Online 官方视频 00:54:23--00:55:39。}

\begin{knowledgebox}{术语消化：输入自由度来自哪里}
\begin{itemize}
  \item \term{type embedding}：告诉模型元素来自 image、radar、map 或 vehicle state；
  \item \term{positional encoding}：注入时间、二维坐标或其他几何关系；
  \item \term{mask/connectivity}：限制哪些元素可以直接通信；
  \item \term{tokenizer/adapter}：把不同原始维度投影到共享 residual width；
  \item \term{inductive bias}：并未消失，而是从 core block 外移到接口、位置与连接规则。
\end{itemize}
\end{knowledgebox}

\teachervoice{Karpathy 的口语总结是“把所有东西切碎，扔进 mix，让 self-attention 决定怎样交流”。这句话抓住工程便利，也容易被误读。真正可复现的版本必须回答：怎么切、怎样标注来源、哪些边允许、用什么目标监督，以及 token 数量是否承受得住。}

\subsection{本章小结}

跨模态迁移的共同模式是把领域对象表示成带类型和位置的元素，再用 attention 路由信息。ViT 仍需要空间位置，Conformer 保留局部卷积，Decision Transformer 依赖轨迹条件，AlphaFold2 使用专门 pair/MSA 表示，异构传感器也要 type embedding 与 connectivity。共享骨架减少融合摩擦，却不替代领域建模。

